% Chapter Template

\chapter{Conclusion and Future Works} % Main chapter title

\label{Chapter 5} % Change X to a consecutive number; for referencing this chapter elsewhere, use \ref{ChapterX}

\lhead{Chapter 5. \emph{Conclusion and Future Works}} % Change X to a consecutive number; this is for the header on each page - perhaps a shortened title

\section{Conclusion}
This work presents a comprehensive evaluation of open-source large language models within an extended multi-agent clinical simulation framework. By moving beyond static QA-style testing and instead modelling doctor–patient interactions, sequential evidence acquisition, and controllable bias conditions, the study provides a more realistic assessment of LLM behaviour in diagnostic settings. \\

The results demonstrate that while medically fine-tuned models exhibit strong domain familiarity, they often underperform compared to instruction-tuned general-purpose models when the task requires multi-step reasoning, decision stability, and structured output adherence. These findings underscore that domain specialization alone is insufficient for reliable diagnostic performance; instead, robust instruction-following ability and strong general reasoning play a critical role in managing the complexities of interactive clinical workflows.  \\

Overall, this project highlights the need for more holistic evaluation frameworks, encourages careful model selection for medical applications, and establishes a foundation for future research on improving the reliability, fairness, and interpretability of open-source LLMs in clinical decision-making.

\section{Future Scope}
The current framework establishes a strong foundation for evaluating LLMs in realistic diagnostic scenarios, but several extensions can further enhance its capability. Future work may focus on integrating multimodal inputs, improving model robustness through targeted fine-tuning, expanding evaluations across diverse patient populations, and developing more quantitative behavioural metrics to better capture reasoning quality and bias patterns.

\begin{itemize}
    \item \textbf{Multimodal integration:} Incorporate radiology images, lab charts, ECGs, and other clinical modalities.
    \item \textbf{Specialized training:} Explore adapters or fine-tuning strategies to improve reasoning stability without catastrophic forgetting.
    \item \textbf{Broader coverage:} Evaluate models across additional languages, demographics, and medical specialties.
    \item \textbf{Safety and uncertainty:} Add calibration tools, hallucination detection, and risk-aware decision modules.
    \item \textbf{Behavioural analytics:} Develop metrics for hallucination rate, test-ordering rationality, and bias sensitivity.

    \item \textbf{Partial-credit evaluation:} Move beyond binary correctness by awarding graded scores (0–1) when the model identifies the correct organ system or disease category but not the exact diagnosis, capturing deeper clinical reasoning quality.
\end{itemize}
